% 本文件由 LaTeX 排版 Agent 根据有效 translation.json 自动生成。
% author=Mar Jacob; work_id=0861; title=关于古里亚与沙穆纳的讲道

\chapter{关于古里亚与沙穆纳的讲道}

% structure: p0001 source=h1 target=omitted reason=page-title-already-rendered-by-parent

\Person{沙穆纳}与\Person{古里亚}，这两位殉道者在苦难中使自己声名显赫，\par

以爱要求我讲述他们光辉的事迹。\par

信德的冠军们，教义呼唤我，\par

让我前去观看他们的竞赛与冠冕。\par

右手之子，与左手争战者，\par

今日召我吟诵他们奇妙争战的故事：\par

朴素的老人，如英雄般投入战斗，\par

在血战中高贵地彰显了自己。\par

他们是吾土之盐，因而土地被调和，\par

其味得以恢复，因不信而变得寡淡。\par

金灯台，充满十字架的油，\par

借此照亮了我们全域，那已转为黑暗之处。\par

两盏灯，当万风肆虐时\par

每一\Emph{种类}错误，其光并未熄灭；\par

良善的工人，从黎明起劳作\par

在天主之家蒙福的葡萄园中，极其忠诚地。\par

吾土之堡垒，成了我们的屏障\par

对抗四周围困我们的各场战争中一切掠夺者。\par

平安的港湾，也是所有困苦者的退避之所，\par

是每个需要援助者的安歇之枕。\par

两颗宝贵的珍珠，它们曾是\par

我主\Person{阿布加尔}（亚兰人之子）新娘的装饰。\par

他们是教师，以血实践其教导，\par

其信德因苦难而显明。\par

在他们身体上，他们书写了天主之子的事迹\par

以密布全身的梳刮与鞭打的\Emph{痕迹}。\par

他们彰显爱，不仅凭口中的言语，\par

而是凭酷刑与肢体的撕裂。\par

为爱天主之子，他们舍弃了身体。\par

因为爱者理应为爱舍弃自己。\par

火与剑验证了他们的爱，\Emph{何其真实}；\par

他们的颈项比地上炉火中\Emph{炼过的}银子更美。\par

他们仰望天主，因见其崇高之美，\par

因此他们轻视为他所受的苦难。\par

正义的太阳在他们心中升起。\par

他们受其光照，以\Emph{他的}光芒驱散了黑暗。\par

他们嘲笑虚妄的偶像，那是谬误引入的，\par

充满光明，因天主之子的信德而生动。\par

主的爱如烈火在他们心中。\par

偶像崇拜的一切荆棘都无法在它面前站立。\par

他们的爱坚定不移地专注于天主：\par

因此他们蔑视那嗜血的刀剑。\par

以淳朴\Emph{兼具}智慧，他们立于审判厅，\par

正如真理之师所命令的。\par

他们鄙视亲友家族，因此淳朴。\par

而且，财产与财富在他们眼中毫无价值。\par

\Emph{不仅淳朴}：在审判厅中，他们也以蛇的智慧\par

谨慎持守天主之家的信德\Emph{也}。\par

当蛇被抓住击打时，它护住头，\par

却将整个身体暴露给捕者。\par

只要头\Emph{免受伤害}，生命就留存于它之内。\par

但若头部被击，生命便成为毁灭的\Emph{猎物}。\par

灵魂的头是人的信德。\par

若此得保全\Emph{无损}，生命也得保全。\par

即使全身因打击而皮开肉绽，\par

\Emph{然而}，只要信德犹存，灵魂就活着。\par

但若信德被不信\Emph{击倒}，\par

灵魂失丧，生命已从那人消逝。\par

\Person{沙穆纳}与\Person{古里亚}，作为信德之人\par

谨慎防备，以免信德被逼迫者\Emph{击倒}。\par

因他们知道，若信德被保存，\par

灵魂与身体都免于毁灭。\par

因此，他们关切自己的信德，\par

以免那隐藏他们生命的东西被\Emph{击倒}。\par

他们交出身体，任人击打与脱臼，\par

甚至每\Emph{一种}酷刑，以使信德不被\Emph{击倒}。\par

又如蛇躲避打击保护头，\par

他们也将信德藏于心中。\par

身体受击，忍受鞭伤与苦痛。\par

但他们心中的信德并未被推翻。\par

口舌说话时，将灵魂出卖给死亡，\par

舌头如刀剑，进行杀戮。\par

从口中生出生命与死亡。\par

否认\Emph{者}死，宣认\Emph{者}活，\Emph{口舌}操此权柄。\par

否认是死亡，宣认中存灵魂生命。\par

口舌如审判官，操二者权柄。\par

口中的言语为死亡开门，\par

同样，它也召唤生命，照亮那人。\par

甚至强盗，因一句信德之言\par

赢得王国，成为充满祝福的乐园的继承人。\par

恶毒的审判官也要求殉道者、右手之子\par

仅凭口舌之词，他们应亵渎。\par

但如同坚守信德的真丈夫，\par

他们未出一言以使不信得逞。\par

\Person{沙穆纳}，我们信德之美，谁能充分述说你？\par

我的口舌为赞美你而过于狭窄，过于卑微，不堪称你。\par

你的真理是你的美，你的冠冕是你的苦难，你的财富是你的鞭痕，\par

因你的打击，你争战之美壮丽非凡。\par

我国以你为荣，如满藏金的宝库。\par

因你于我们是财富，是无可窃夺的珍藏。\par

\Person{古里亚}，殉道者，我们信德的坚定英雄，\par

谁能足够述说你的神圣之美？\par

看！你身上的酷刑如绿柱石宝石镶嵌，\par

你颈上的刀剑如精金链。\par

你身上的鲜血是荣耀之袍，满溢美丽，\par

你背上的鞭打是太阳无法比拟的衣裳。\par

因这些丰盛的苦难，你光彩照人、俊美无比。\par

你的美辉煌灿烂，因你承受的剧痛\Emph{如此}。\par

\Person{沙穆纳}，我们的财富，你比富人更富。\par

因为看！富人立于你的门口，求你周济。\par

你的村庄渺小，你的国家贫瘠：那么，是谁赐予你\par

使乡村与城市的主人都来讨好你？\par

看！身着长袍与法衣的审判官们\par

从你的门槛取些尘土，\Emph{仿佛}那是生命的良药。\par

十字架是富有的，且为朝拜它的人增添财富；\par

而它的贫穷藐视世上的一切财富。\par

沙姆纳和古里亚，穷人的儿子，看哪！你们门前\par

富人俯首，好从你们那里领受\Emph{所需}的供应。\par

天主子在贫穷与匮乏中\par

向世界显示，世上一切财富皆虚无，\par

\Emph{祂的门徒}——全是渔夫，全是穷人，全是弱者，\par

全是无名之辈——因祂的信德而变得显赫。\par

一位渔夫，其\Quote{村庄}也是渔夫之家，\par

祂立他为十二人之首，是的，为家之首领。\par

一位制帐篷者，从前曾是迫害者，\par

祂攫住他，使他成为信德的蒙选器皿。\par

沙姆纳和古里亚来自并不富足的村庄，\par

看哪！在一座大城中他们成了领主；\par

而城中显要、审判官也站在他们门前，\par

并恳求他们的施舍，以满足其需要。\par

因他们宣认了天主子信德的宣认，\par

这些蒙福之人获得了无法计算的财富。\par

祂自己成了贫穷的，并使贫穷的成为富足；\par

看哪！藉着祂的贫穷，整个受造界得到了富足。\par

蒙选的殉道者与谬误作战，\par

并在宣认天主子的事上，如勇士般屹立不动。\par

他们进去，在审判官面前以无畏的神色宣认祂，\par

好使祂也宣认他们，正如他们在祂的圣父前宣认了祂。\par

外邦人的战争如暴风般兴起攻击他们；\par

但十字架是他们的舵手，引领他们前行。\par

他们被要求向无生命的偶像祭献，\par

但他们没有偏离对天主子的宣认。\par

偶像崇拜的风迎面吹来，\par

但他们自己如同堆积的磐石，抗击飓风。\par

谬误如迅疾的旋风攫取他们；\par

但因他们受十字架刑罚的庇护，未能伤害他们。\par

恶者放出所有犬吠叫，要咬他们；\par

但因他们以十字架为杖，尽都驱散了。\par

但谁能足够述说他们的竞赛，\par

或他们的苦难，或他们肢体的撕裂？\par

或谁能描绘他们加冕的画面，\par

他们如何从竞赛中荣耀地升上？\par

他们进入审判，但毫不介意审判官；\par

也不焦虑受审时该说什么。\par

审判官威吓\Emph{他们}，倍增恐吓之词；\par

并历数酷刑和\Emph{种种}折磨，想惊吓他们。\par

他说大话，想藉恐惧和威吓，\par

也藉恫吓，使他们倾向于祭献。\par

然而斗士们藐视恫吓、威吓、\par

审判的判决，以及所有身体的死亡；\par

他们预备自己忍受侮辱、鞭打、打击、\par

挑衅、被拖曳、被焚烧；\par

以及监禁、锁链，和一切邪恶之事，\par

并一切酷刑和一切苦难，且始终喜乐。\par

他们不惊恐、不惧怕、不沮丧，\par

酷刑的锋利也未使他们屈从祭献。\par

他们鄙视身体，视如地上的粪土：\par

因为他们知道，越被击打，其美丽越增；\par

审判官越增加恫吓来惊吓他们，\par

他们越显出对他的轻蔑，毫不畏惧他的威胁。\par

他不断告诉他们，他为他们准备了什么酷刑；\par

他们也不断告诉他，为他保留的地狱。\par

藉他所告诉\Emph{他们}的事，他试图吓唬他们去祭献；\par

他们则对他讲论那边可畏的审判。\par

真理比智慧的言语更智慧，\par

而谎言，纵使极力粉饰，也极其可憎。\par

沙姆纳和古里亚不断讲论真理，\par

而审判官继续吐出谎言。\par

因此他们不惧怕他的威胁，\par

因为他反对真理的一切恫吓，被\Emph{他们}视为空响。\par

他们鄙视、轻蔑、嘲笑世上的交往，\Emph{是的}，舍弃了；\par

并且无意返回，或\Emph{再}进入其中。\par

他们从审判之地定睛要离去，\par

前往他们众人相聚之处，即新世界的生命。\par

他们既不关心财产，也不关心房屋，\par

也不关心这个充满邪恶的世界的利益。\par

在光明世界中，他们的心与天主一同被吸引，\par

并定意往\Quote{那}地方去；\par

他们仰望刀剑，指望它来作一座桥，\par

好让他们渡到他们所渴慕的天主那里。\par

他们视这世界如同一座小帐篷，\par

而那边则是一座充满美景的城；\par

他们急着藉刀剑离开此地，\par

前往那光明之乡，那里充满对配得之人的祝福。\par

审判官下令悬起他们的手臂，\par

他们毫不留情地拉开他们，使其处于痛苦煎熬中。\par

魔鬼的狂怒将怒气吹入审判官的心，\par

并使他苦待那坚定者，\Emph{煽动他}粉碎他们；\par

他在高处与深处之间拉开他们以折磨他们：\par

两边都惊奇，\Emph{见}他们忍受了多少。\par

天地都惊异于老者们的躯体，\par

\Emph{见}它忍受了多少痛苦，却在\Emph{其}苦难中不呼救。\par

他们衰弱的身体被吊起、拖曳着手臂，\par

然而一片寂静，无人呼救或抱怨。\par

凡看见他们竞赛的人都惊讶，\par

\Emph{见}那伸展的身形如何\Emph{平静地}承受\Emph{加于其上}的折磨。\par

撒殚也惊讶于他们无暇的身躯，\par

\Emph{见}他们承受何等沉重的苦难，毫无呻吟。\par

是的，天使们也因\Emph{他们的}刚毅而欢喜，\par

\Emph{见}它如何耐心地承受那如此可怕的竞赛。\par

但如同等待冠冕的战士，\par

他们心中毫无倦意。\par

不，倒是审判官倦了；是的，他惊异：\par

但高尚者\Emph{在他面前}在苦难中毫无倦意。\par

他问他们是否同意祭献；\par

但口因痛苦而无法言语。\par

于是迫害者增加折磨，\par

直到不给言语留余地。\par

因肢体所受的折磨，口沉默不语；\par

然而，他的意志却如同英雄一般，从自身获得了坚韧的力量。\par

迫害者有祸了！他们何其缺乏正义！\par

但光明之子——他们如何披戴了信德！\par

他们要求发言，却无发言之处，\par

因为口舌之言已被痛苦禁止。\par

身体被紧紧束缚，口舌沉默，无法\par

在不义受审时说出话语。\par

殉道者又能做什么呢？他无力说出——\par

当被问及他是否不愿献祭时——\par

满有信德的老人们全都沉默，\par

因痛苦而无法说话。\par

然而他们仍被审问：若人沉默不语，\par

被问时，又怎能赞同所言之语？\par

但老人们，为避免被认为赞同，\par

用记号清楚地表达了他们理当说出的话。\par

他们摇头，代替言语，以无声的记号显示\par

内在新人的决心。\par

他们垂头，在痛苦中表示\par

他们不准备献祭，人人都懂其意。\par

当他们尚能说话时，便以言语宣认；\par

但当痛苦禁止言语时，他们便以无声的记号发言。\par

他们论信德，既有声音，亦无声响：\par

以致无论说话还是沉默，他们都\Emph{同样}坚定不移。\par

谁能不对生命之路惊叹：它多么狭窄，\par

对愿行其上者又多么笔直？\par

谁能不惊异\Emph{看到}：当意志警醒而备妥时，\par

它便极其宽广，满有光明，给行在其中者？\par

路旁有沟渠；也满是陷阱；\par

若稍偏离，沟渠便接纳了他。\par

仅在那无声的记号之间，左右之隔，\par

在\Quote{是}与\Quote{非}之上，罪与义得以成立。\par

仅凭一个无声的记号，蒙福者们清楚地表明他们不献祭，\par

因这一个无声的记号，这路便引他们进入伊甸园；\par

而若这同一无声的记号稍作倾斜、向下偏转\par

朝向深渊，老人们的路便会通向地狱。\par

他们向上做记号，\Emph{表示}他们预备向上攀登；\par

因那记号，他们上升并与天上者相合。\par

在记号与记号之间，是乐园与地狱：\par

他们做了不献祭的记号，便承受了天国。\par

即使他们沉默时，仍是天主之子的辩护者：\par

因为信德不在于言语之多。\par

他们的刚毅是洪亮的宣认，\par

仿佛张口以记号宣示了\Emph{他们的}信德；\par

人人都知道他们虽沉默，却说什么，\par

天主之家的信德被丰富和增长；\par

谬误因两位老人而蒙羞——他们虽不言语，\par

却战胜了它；他们保持沉默，信德却屹立不摇。\par

尽管法官口出暴烈言辞，\par

皇帝的命令可怕且狂暴，\par

异教厚颜张口，\par

高声喧嚷，老人们却因痛苦而沉默，\par

\Emph{然而}命令落空，法官的声音被淹没，\par

未发一言，殉道者无声的记号却夺得了棕枝。\par

说话、喧嚷与鞭打之声在左；\par

深沉沉默与受苦站立在右；\par

凭老人们举过头顶的一个无声记号，\par

信德之首被高举，谬误蒙羞。\par

说话者落败，沉默者得胜：\par

因他们无语，却以记号发出了信德之言。\par

他们被带下，因沉默而获胜；\par

他们被捆绑，受威胁\Emph{仍要}制胜他们。\par

捆绑和黑暗无光的监牢，在殉道者看来\par

不足挂齿——\Emph{反而}如同无尽的光明。\par

\Emph{处于}无饼、无水、无光之境，\par

他们仍感欣慰，因天主之子的爱。\par

法官下令将他们双脚吊起\par

头朝下，全不义的判决：\par

沙穆纳被吊起，头朝下；他祈祷，\par

在痛苦中纯洁而滤净的\Emph{祈祷}。\par

甘美果实挂在审判厅的树上，\par

其滋味与香气使天上的居民也惊奇。\par

他的身体受苦，但信德健全；\par

他本人被紧绑，但祈祷因他的行为而无羁绊。\par

因为祈祷丝毫不被任何事所转，\par

也\Emph{无物}可阻挡——哪怕刀剑，哪怕火焰。\par

他的形体颠倒，但\Emph{他的}祈祷无拘无束，\par

其路径笔直向上，直达天使的居所。\par

被选的殉道者的苦楚愈增，\par

从他口中愈发出全部宣认。\par

殉道者们渴望磨利的剑，满怀深情，\par

视之如充满财富的宝藏。\par

天主之子在世成就了一个新工作——\par

可怖的死亡竟为多人所渴求。\par

人迎向刀剑，是闻所未闻之事，\par

除非他们是耶稣以祂的十字架召募来事奉祂的人。\par

死亡是苦涩的，人人皆知！看，从太初以来：\par

唯独殉道者不以为被杀戮是苦涩的。\par

他们看见磨利的刀剑，便嘲弄它，\par

以微笑迎接：因为那是他们得冠冕的机缘。\par

仿佛它是什么可憎之物，他们任身体被鞭打：\par

即使爱惜它，也不将它从痛苦中收回。\par

他们等候刀剑，刀剑便赐下桂冠：\par

因为他们期盼它，它便来迎接他们，正如他们所愿。\par

天主之子以祂的十字架击杀死亡；\par

正因死亡已被击杀，它便不再使殉道者受苦。\par

受伤之蛇，人可无惧把玩；\par

已死的狮子，连懦夫也能拖走：\par

我们主的大蛇被祂的十字架击碎；\par

可怖的狮子被天主之子以祂的苦难击杀。\par

祂将死亡牢牢捆绑，将之推倒，践踏在阴府门口；\par

\Emph{如今}凡愿意的，都就近嘲弄它，因它已被杀。\par

这些老人——沙穆纳和古里亚——嘲弄死亡，\par

如同嘲弄那被天主之子所杀的狮子一样。\par

那曾在树间击杀亚当的大蛇，\par

\Emph{谁}能抓住他，只要他不喝那十字架的血？\par

天主圣子以祂的十字架刑粉碎了巨龙；\par

看哪！童稚与老翁都嘲弄那受伤的蛇。\par

那狮子被刺伤了，被那\Emph{刺穿}天主圣子肋膀的长矛；\par

凡愿意的人都能践踏他，甚至嘲弄他。\par

天主圣子——祂是一切美善的根源，\par

每个口唇都当颂扬祂。\par

祂亲自以从伤口流出的血迎娶了新娘，\par

祂向婚宴的宾客要求他们颈项的血作为婚嫁之礼。\par

婚宴的主赤身悬挂在十字架上，\par

凡来赴宴的，祂都让祂的血滴落在他身上。\par

沙穆纳和古里亚为祂的缘故交出了自己的身躯，\par

去忍受痛苦、酷刑与各种形式的灾祸。\par

他们仰望祂，正当祂被恶人嘲弄，\par

他们自己也如此忍受嘲弄而无一声呻吟。\par

\Place{埃德萨}因你们的被杀而富足，哦有福的人们：\par

因为你们以你们的冠冕和苦难装饰了她。\par

你们是她的美丽，她的堡垒，她的盐，\par

她的财富和库藏，是的，她的夸耀和一切珍宝。\par

你们是忠信的管家：\par

因为藉着你们的苦难，你们装扮了新娘，使她美丽。\par

那\Place{帕提亚}的女子，她已许配给十字架，\par

以\Emph{你们}为夸耀：因为藉着你们的教导，看哪！她得了光照。\par

你们是她的辩护者；是文士，虽沉默却击败了\par

一切谬误，当其声音在不信中高扬之时。\par

希伯来女儿的那些长老是\Person{彼列}之子，\par

假见证人，他们杀害了\Person{纳波特}，自充为\Emph{真}。\par

\Place{埃德萨}以其两位满有美善的长老胜过她，\par

他们是为天主圣子作证的人，像\Person{纳波特}一样死去。\par

那里有两位，这里也有两位，都是长老；\par

这些被称为证人，那些也是证人。\par

现在让我们看看他们中哪一些是天主拣选的证人，\par

哪一座城因其长老和尊贵之人而蒙爱。\par

看哪！杀害\Person{纳波特}的\Person{彼列}之子是证人；\par

而这里的\Person{沙穆纳}和\Person{古里亚}，同样也是证人。\par

现在让我们看看哪些证人，哪些长老，\par

以及哪一座城能坦然无惧地站在天主面前。\par

那淫妇的证人是\Person{彼列}之子，\par

看哪！他们的羞耻尽都刻画在他们的名字里。\par

\Place{埃德萨}的公义正直的长老，她的证人，\par

如同\Person{纳波特}，他也曾为正义而被杀。\par

他们不像那两个说谎的\Person{彼列}之子，\par

\Place{埃德萨}也不像\Place{熙雍}，她也曾钉死\Emph{主}。\par

她的长老像她自己一样虚假，竟敢\par

邪恶地将无辜的血倾倒在地。\par

\Emph{然而}藉着这里的这些证人，看哪！真理被宣告了。——\par

愿那赐给我们他们冠冕宝库者受赞美！\par

\Emph{这里}结束\WorkTitle{古里亚与沙穆纳讲道词}。\par

\SourceRecord{Translated by B.P. Pratten.；Ante-Nicene Fathers；Vol. 8.；Edited by Alexander Roberts, James Donaldson, and A. Cleveland Coxe.；Buffalo, NY: Christian Literature Publishing Co.,；1886.；<http://www.newadvent.org/fathers/0861.htm>.}
